\documentclass[lettersize,journal]{IEEEtran}

% ============================================================
% PACKAGES
% ============================================================
\usepackage{amsmath,amssymb,amsfonts,amsthm}
\usepackage{algorithm}
\usepackage{algpseudocode}
\usepackage{graphicx}
\usepackage{cite}
\usepackage{url}
\usepackage{xcolor}
\usepackage{booktabs}
\usepackage{multirow}
\usepackage{bm}
\usepackage[caption=false,font=footnotesize]{subfig}
\usepackage{balance}
\usepackage{tikz}
\usetikzlibrary{arrows.meta,positioning,calc,decorations.pathmorphing,
  shapes.geometric,fit,backgrounds,patterns,fadings,shadows.blur}
\usepackage{pgfplots}
\pgfplotsset{compat=1.18}

\hyphenation{op-tical net-works semi-conduc-tor Kolmo-go-rov}

% ============================================================
% THEOREM ENVIRONMENTS
% ============================================================
\newtheorem{corollary}{Corollary}
\newtheorem{lemma}{Lemma}
\newtheorem{remark}{Remark}
\newtheorem{definition}{Definition}
\newtheorem{proposition}{Proposition}
\newtheorem{assumption}{Assumption}

% ============================================================
% CUSTOM COMMANDS
% ============================================================
\DeclareMathOperator{\erfc}{erfc}
\DeclareMathOperator{\diag}{diag}
\DeclareMathOperator*{\argmin}{arg\,min}
\DeclareMathOperator*{\argmax}{arg\,max}
\newcommand{\vect}[1]{\mathbf{#1}}
\newcommand{\mat}[1]{\mathbf{#1}}
\newcommand{\set}[1]{\mathcal{#1}}
\newcommand{\norm}[1]{\left\|#1\right\|}
\newcommand{\abs}[1]{\left|#1\right|}
\newcommand{\R}{\mathbb{R}}
\newcommand{\E}{\mathbb{E}}

% ============================================================
\usepackage[active,tightpage]{preview}
\setlength\PreviewBorder{0.6pt}
\begin{document}
\begin{preview}
\begin{tikzpicture}
\begin{axis}[
  width=\columnwidth, height=4.5cm,
  xlabel={Time (s)}, ylabel={Normalized concentration},
  xmin=0, xmax=30,
  ymin=-0.05, ymax=1.1,
  grid=major, grid style={gray!20},
  legend style={font=\scriptsize, at={(0.97,0.97)}, anchor=north east},
  tick label style={font=\scriptsize},
  label style={font=\small},
]
% Simulated mean received concentration (CIR shape)
% Real 40-trial testbed data (data/MacroscaleTestbed/dataset_SISO_testbed.csv); band = mean $\pm$ s.d.
\addplot[blue!25, draw=none, fill opacity=0.4, forget plot]
  coordinates {(0.0,0.001) (0.5,0.002) (1.0,0.074) (1.5,0.583) (2.0,0.863) (2.5,0.992) (3.0,1.002) (3.5,0.967) (4.0,0.897) (4.5,0.812) (5.0,0.733) (5.5,0.658) (6.0,0.597) (6.5,0.544) (7.0,0.497) (7.5,0.461) (8.0,0.427) (9.0,0.364) (10.5,0.293) (12.0,0.249) (13.5,0.200) (15.0,0.169) (16.5,0.143) (18.0,0.122) (19.5,0.104) (21.0,0.090) (22.5,0.080) (24.0,0.072) (25.5,0.065) (27.0,0.059) (28.5,0.054) (30.0,0.051) (30.0,0.024) (28.5,0.026) (27.0,0.030) (25.5,0.034) (24.0,0.038) (22.5,0.044) (21.0,0.052) (19.5,0.062) (18.0,0.075) (16.5,0.092) (15.0,0.112) (13.5,0.138) (12.0,0.169) (10.5,0.214) (9.0,0.277) (8.0,0.334) (7.5,0.369) (7.0,0.409) (6.5,0.456) (6.0,0.505) (5.5,0.568) (5.0,0.646) (4.5,0.724) (4.0,0.813) (3.5,0.905) (3.0,0.978) (2.5,0.902) (2.0,0.672) (1.5,0.289) (1.0,0.000) (0.5,0.000) (0.0,0.000)} -- cycle;
\addplot[blue, very thick, no markers]
  coordinates {(0.0,0.000) (0.5,0.000) (1.0,0.024) (1.5,0.436) (2.0,0.767) (2.5,0.947) (3.0,0.990) (3.5,0.936) (4.0,0.855) (4.5,0.768) (5.0,0.690) (5.5,0.613) (6.0,0.551) (6.5,0.500) (7.0,0.453) (7.5,0.415) (8.0,0.380) (9.0,0.321) (10.5,0.254) (12.0,0.209) (13.5,0.169) (15.0,0.141) (16.5,0.117) (18.0,0.099) (19.5,0.083) (21.0,0.071) (22.5,0.062) (24.0,0.055) (25.5,0.050) (27.0,0.044) (28.5,0.040) (30.0,0.038)};
\addlegendentry{Mean (40 trials)}
\addplot[red, dashed, very thick, no markers]
  coordinates {(0.0,0.000) (0.5,0.018) (1.0,0.308) (1.5,0.614) (2.0,0.764) (2.5,0.807) (3.0,0.796) (3.5,0.760) (4.0,0.715) (4.5,0.668) (5.0,0.622) (5.5,0.579) (6.0,0.539) (6.5,0.502) (7.0,0.469) (7.5,0.439) (8.0,0.412) (9.0,0.364) (10.5,0.307) (12.0,0.263) (13.5,0.229) (15.0,0.201) (16.5,0.179) (18.0,0.160) (19.5,0.144) (21.0,0.131) (22.5,0.119) (24.0,0.110) (25.5,0.101) (27.0,0.094) (28.5,0.087) (30.0,0.081)};
\addlegendentry{CIR fit ($R^2{=}0.913$)}
\end{axis}
\end{tikzpicture}
\end{preview}
\end{document}
